% SHARD-ID: RC-06D-RING-NORM-ELLIPTIC
% SHARD-TITLE: Curve Counts as Twisted Ring Norms II: the Ordinary Elliptic Curve
% SHARD-SUMMARY: The natural tensor of an ordinary elliptic curve is its lifted Frobenius acting on the flat torus: the count is a Lefschetz number, the grading is form degree, the ket/bra doubling is the Hodge decomposition, the Polya--Hilbert unitary is the pullback on harmonic one-forms, the Ramanujan property is the degree of the lift, and the gauge freedom is the lattice basis whose classes are ideal classes.
% SHARD-SUMMARY: The Fourier pullback is an injection whose only finite orbit is the zero mode; digit expansions in base pi have bounded cyclic carries yet no nonnegative automaton counts the points.
% SHARD-KEYWORDS: ordinary elliptic curve, Deuring lift, toral endomorphism, Lefschetz number, Hodge decomposition, ring norm, supertrace, Polya-Hilbert, Hasse bound, gauge, ideal class, Fourier shift, digit expansion, carry automaton
% SHARD-NOTES: none
% SHARD-SCRIPTS: scripts/ring_norm_tensor.py

\section{Curve Counts as Twisted Ring Norms II: the Ordinary Elliptic Curve}
\label{sec:ring-norm-genus-one-tier-a}

Tier A of the campaign of \cref{sec:ring-norm-genus-one}; notation of
\cref{def:lifted-frobenius-torus,def:elliptic-norm-tensor}; proofs in
\texttt{notes/ring-norm-tensor/astra-proofs.md} T1--T2.

\subsection{Tier A: the ordinary elliptic curve}

\begin{theorem}[Frobenius on the lifted torus; the count is a ring norm]
\label{thm:toral-frobenius-count}
\claimstatus{thm:toral-frobenius-count}{proved-conditional}
With $M$ the integral matrix of \cref{def:lifted-frobenius-torus}, $\chi_M(X) = X^2-aX+\qs$
and for every $n\ge1$ the toral endomorphism $M^n$ has exactly $\det(1-M^n) = |1-\pi^n|^2 =
N_{K/\mathbb Q}(1-\pi^n) = \Ncount{n}$ fixed points, all nondegenerate of index $+1$, so that
$\Ncount{n} = 1-\Tr M^n+\det M^n$ is the Lefschetz number; moreover $\Ncount{n} = |O/(\pi^n-1)O|$,
the index of the principal ideal $(1-\pi^n)$ in the ring $O$, which under
\cref{cit:lenstra-structure} is the point group $E(\Fqn{n})$ itself as an $O$-module.
\end{theorem}

\begin{proof}
Sketch (T1.1). $\pi$ preserves $\Lambda$, so $M$ is integral with eigenvalues $\pi,\bar\pi$, % bare-ok
determinant $|\pi|^2 = \qs$ by \cref{cit:degree-norm-index,cit:hasse-degree-bound} and trace $a$
by \cref{cit:frob-count-formula}. The fixed points of $M^n$ form $\mathbb Z^2/(M^n-1)\mathbb Z^2$ of
order $|\det(M^n-1)|$; $\det(1-M^n) = (1-\pi^n)(1-\bar\pi^n)>0$ gives the index. Lefschetz
(\cref{cit:toral-lefschetz}) with the actions $1,M^T,\qs$ in degrees $0,1,2$ gives the count.
\end{proof}

\begin{proposition}[Dynamical zeta equals Hasse--Weil zeta]
\label{prop:toral-zeta}
\claimstatus{prop:toral-zeta}{proved-conditional}
$\exp\bigl(\sum_n\#\operatorname{Fix}(M^n)\uz^n/n\bigr) = (1-a\uz+\qs\uz^2)/((1-\uz)(1-\qs\uz))
= Z(E,\uz)$; by M\"obius inversion the primitive periodic orbits of $M$ on $\mathbb T^2$ of each
length are equinumerous with the closed points of $E$ of each degree, so the gas side and the
transfer side are both carried by the dynamical system $(\mathbb T^2,M)$; the bijection is
noncanonical.
\end{proposition}

\begin{theorem}[Hodge doubling and the elliptic norm tensor]
\label{thm:hodge-doubling-tensor}
\claimstatus{thm:hodge-doubling-tensor}{proved-conditional}
On $\mathbb C/\Lambda$, $M^*dz = \pi dz$ and $M^*d\bar z = \bar\pi d\bar z$. The double of the % bare-ok
Hodge ket bond is $H^*(\mathbb T^2,\mathbb C)$ graded by form degree, with transfer
$\Edb = \operatorname{diag}(1,\bar\pi,\pi,\qs)$ in the basis $(1,d\bar z,dz,dz\wedge d\bar z)$ and
$\Gdb = \operatorname{diag}(1,-1,-1,1)$: $(+,+) = H^0$, the mixed lines are $H^{0,1},H^{1,0}$, and
$(-,-) = H^2$ with eigenvalue $\pi\bar\pi = \qs$, the degree. The elliptic norm tensor
(\cref{def:elliptic-norm-tensor}) is all-even, $\Psi_P(n) = (1-\pi^n)\phi^{\otimes n}$,
$\|\Psi_P(n)\|^2 = \Ncount{n}$, the untwisted norm is $|1+\pi^n|^2$, and the ring zeta is
$Z(E,\uz)$. The fermionic modes are the odd-degree forms; the physical letters are bosonic;
the physical state has Schmidt rank one and its configurations do not enumerate points.
\end{theorem}

\begin{proposition}[Fourier pullback is an injection with one finite orbit; the flat supertrace]
\label{prop:fourier-shift-flat-trace}
\claimstatus{prop:fourier-shift-flat-trace}{proved-conditional}
The pullback $U = M^*$ on $L^2(\mathbb T^2)\otimes\Lambda^*(\mathbb C^2)$ sends $e_k\otimes\omega$ to % bare-ok
$e_{M^Tk}\otimes\Lambda^*(M^T)\omega$; $k\mapsto M^Tk$ is injective of index $\qs$ on $\mathbb Z^2$ % bare-ok
(not a permutation), its only finite orbit is $\{0\}$, and the Fourier-diagonal flat supertrace
$\sum_{M^{nT}k = k}\str\Lambda^*(M^{nT}) = \det(1-M^{nT}) = \Ncount{n}$ is contributed by $k = 0$ % bare-ok
alone. The infinite-label tensor $\mathsf A_k = |M^Tk\rangle\langle k|\otimes\operatorname{diag}(1,\pi)$
has $P$-closed ring vector $(1-\pi^n)|0,\dots,0\rangle$: an honest but redundant physical
index; every closed ring except the vacuum mode vanishes.
\end{proposition}

\begin{theorem}[The Polya--Hilbert unitary and the ordinary Hasse bound]
\label{thm:elliptic-ph-unitary}
\claimstatus{thm:elliptic-ph-unitary}{proved-conditional}
On the flat torus the Hodge (L$^2$) metric on harmonic one-forms is scaled by $|\pi|^2$ under
$M^*$ on both $H^{1,0}$ and $H^{0,1}$, so $U_{\mathrm{PH}} = \qs^{-1/2}M^*|_{H^1}$ is unitary iff
$|\pi|^2 = \qs$, which is derived, not assumed: the lifted Frobenius is a conformal similarity
whose oriented area ratio is its covering degree $[\Lambda:\pi\Lambda] = |\pi|^2$ (\cref{cit:degree-norm-index}), % bare-ok
and degree is preserved under good reduction with $\deg\Frob = \qs$ (the clause $d = \deg(\pi) = \qs$ of
\cref{cit:hasse-degree-bound}, not its inequality; review m1).
Hence the two nontrivial eigenvalues of $\Edb$ have modulus $\sqrt{\qs}$ (the Ramanujan
property here) and $a^2<4\qs$, from the positive definite degree form $\deg(m+n\pi) =
m^2+amn+\qs n^2$; no root-size input is used. The real PH model $K\otimes_{\mathbb Q}\mathbb R\cong\mathbb C$
with $\operatorname{Re}(x\bar y)$ (the Rosati trace form up to the factor $2$) is the real form of
$H^1(\mathbb T^2,\mathbb R)$, not the complex $H^1$; a self-adjoint logarithm needs a phase branch
the curve does not specify.
\end{theorem}

\begin{proposition}[Integral markings versus MPS gauge]
\label{prop:integral-marking-gauge}
\claimstatus{prop:integral-marking-gauge}{proved-conditional}
The integral matrix $M$ changes by $\mathrm{GL}_2(\mathbb Z)$ conjugacy under a change of lattice
basis; its class is an ideal-lattice class of $\mathbb Z[\pi]$ (\cref{cit:latimer-macduffee}, all
full ideal lattices, not only invertible ones) whose multiplier order is the endomorphism
order, and the ordinary classification is recovered as Waterhouse's torsor (indexed by
$j$-invariants, up to quadratic twist; review m4). Complex spectral
data forget these classes: all matrices in the isogeny class are $\mathrm{GL}_2(\mathbb C)$-conjugate,
the cohomology matrix is $M^T$ on the dual lattice, and the unitary gauge is the Hodge basis.
The even ket gauge $G\in\mathrm{GL}(1|1)$ acts on $\Edb$ by $G\otimes\bar G$ and fixes all ring
norms; exchanging the two Hodge embeddings is not an even gauge.
\end{proposition}

\subsection{Digits in base $\pi$}

\begin{proposition}[Ordinary digit expansions, cyclic fibres, bounded carries]
\label{prop:cyclic-digits-carries}
\claimstatus{prop:cyclic-digits-carries}{proved}
With $D$ a residue system of $O/\pi O$ ($|D| = \qs$), every class of $O/\pi^nO$ has a unique
expansion $\sum_{i<n}d_i\pi^i$. The cyclic map $D^n\to O/(\pi^n-1)O$ need not be onto (for
$O = \mathbb Z[i]$, $\pi = 1+2i$, $D = \{0,\pm1,\pm2\}$ the target $O/(2i)$ has four classes and the
digits hit two; $|1-\pi^2|^2 = 32>25$ excludes surjectivity at $n = 2$ for any five digits; this
is $y^2 = x^3+x$ over $\mathbb F_5$). Two words have the same image iff they admit integral cyclic
carries $d_i-e_i = \pi c_{i+1}-c_i$, $c_n = c_0$, which are unique and bounded by
$|c_i|\le\max_{d,e}|d-e|/(\sqrt{\qs}-1)$: the cyclic equivalence relation has a finite carry
automaton.
\end{proposition}

\begin{proposition}[What finite carries do not count]
\label{prop:carry-automaton-no-count}
\claimstatus{prop:carry-automaton-no-count}{proved}
No finite weighted adjacency matrix and no trace-class operator has power traces
$\Ncount{n}$ (\cref{thm:no-ungraded-trace}), so no automaton with nonnegative transition counts
counts the classes of $O/(\pi^n-1)$ by its unweighted rings; the carry automaton recognises
equality of represented words and does not count classes with weight one. A graded signed
transfer ($\operatorname{diag}(1,\qs)$ even, $M$ odd) has supertrace $\Ncount{n}$ but is not by itself a
doubled-Kraus factorisation; the elliptic norm tensor supplies that. The drafted ``carries must
be unbounded'' was wrong: cyclic carries are bounded here.
\end{proposition}

